\documentclass[11pt,a4paper]{article}
\usepackage[margin=25mm]{geometry}
\usepackage{amsmath,amssymb,amsthm,hyperref}
\newtheorem{lemma}{Lemma}
\newcommand{\gen}[1]{\langle#1\rangle}
\title{Alternating attachment and a central character}
\author{Human Compression project: bounded L044 proof component}
\date{}
\begin{document}\maketitle
These are explicit applications of alternating generation and the
central-character interface used in the accepted common theory.
They recover the two borrowed interfaces quoted as Propositions~5.8
and~5.9 of Cameron--Fernandes--Leemans--Mixer,
\emph{Highest rank of a polytope for $A_n$},
arXiv:1605.09173v1, physical pages~25--26. L044 attributes both to Fernandes--Leemans--Mixer,
\emph{All alternating groups $A_n$ with $n\geq12$ have polytopes
of rank $\lfloor(n-1)/2\rfloor$}, SIAM J. Discrete Math.~26 (2012),
Proposition~3.3 and Proposition~5.4. We retain that attribution;
this note does not claim a review of that separate original.
There is no new classification or formalization claim.

\begin{lemma}[Alternating support gluing]
For finite sets $Y,Z$ with $|Y|,|Z|\geq3$ and $|Y\cap Z|\geq2$,
\[
                    \gen{A_Y,A_Z}=A_{Y\cup Z}.
\]
\end{lemma}
\begin{proof}
Choose $p,q\in Y\cap Z$. A three-cycle $(p\ x\ y)$ lies in one
of the two alternating groups if $x,y$ belong to a common set.
For the remaining case $x\in Y\setminus Z$, $y\in Z\setminus Y$,
right-action multiplication gives
\[
                (p\ q\ x)^{-1}(p\ q\ y)=(p\ x\ y).
\]
Thus all three-cycles through $p$ belong to the generated group.
They generate the alternating group: every three-cycle off $p$
is a product of two such cycles, and every even permutation is a
product of three-cycles. The reverse containment is immediate.
\end{proof}

\begin{lemma}[Four-point alternating attachment]
Let $H$ be any transitive permutation group on $X$, $|X|\geq5$,
fixing four new points $u,v,w,z$. Choose $i\in X$ and put
\[
               a=(i\ u)(v\ w),\qquad b=(u\ v)(w\ z).
\]
Then $K=\gen{H,a,b}$ is $A_{X\cup\{u,v,w,z\}}$ if $H$ is even,
and the full symmetric group otherwise.
\end{lemma}
\begin{proof}
For every $j\in X$, transitivity supplies an $H$-conjugate
$a_j=(j\ u)(v\ w)$. Products $a_j a_l$ for distinct $j,l\in X$
supply all three-cycles
through $u$ on $X\cup\{u\}$. Hence $K$ contains
$A_{X\cup\{u\}}$. Conjugating by $b$ also gives
$A_{X\cup\{v\}}$. Lemma~1 combines these to give
$A_{X\cup\{u,v\}}$.
Conjugating this subgroup by $a$ gives
$A_{X\cup\{u,w\}}$, since $a$ exchanges $i,u$ within the old
support and takes $v$ to $w$. Gluing yields
$A_{X\cup\{u,v,w\}}$. Finally its conjugate under $b$ is
$A_{X\cup\{u,v,z\}}$, and a last gluing yields the full
alternating group on all points. Both $a,b$ are even. Thus
$K$ equals this alternating group when $H$ is even; if $H$
contains an odd permutation, $K$ contains the other coset as well
and is the full symmetric group.
\end{proof}

\begin{lemma}[One central twist, with marked scope]
Let $H=\gen{\alpha_0,\ldots,\alpha_{d-1}}$ be a string group
generated by involutions, and let $\tau$ be an involution commuting
with $H$ but not belonging to $H$. Fix $k$ and put
\[
 H^*=\gen{\alpha_j\tau^{\delta_{jk}}:0\leq j<d}.
\]
Then $H^*$ is abstractly $H$ or $H\times C_2$.
If $\tau\notin H^*$, projection to $H$ is a marked isomorphism;
in particular the original marking has IP if and only if the
twisted marking has IP.
\end{lemma}
\begin{proof}
The hypotheses identify $\gen{H,\tau}$ with $H\times\gen\tau$.
Projection restricts to a surjection $H^*\to H$, with kernel
either trivial or $\gen\tau$. In the latter case, removing
$\tau$ from the one twisted generator shows that $H\leq H^*$,
so $H^*=H\times\gen\tau$. In the former case the projection
is an isomorphism taking each distinguished twisted generator
to the corresponding original generator. It therefore takes
every marked parabolic, and every intersection of parabolics,
isomorphically to its original counterpart.

Equivalently, the trivial-kernel case is the graph of a
character $\chi:H\to C_2$ taking $\alpha_k$ to $1$ and the
other distinguished generators to $0$. This is the direct
instantiation of the accepted central-character method.
The product case makes no additional claim about marked IP;
an abstract product decomposition alone would not prove it.
\end{proof}

\paragraph{Exact source boundary.}
Lemma~2 proves the complete group-generation conclusion of source
Proposition~5.8; it does not assert that adjoining the two involutions
preserves IP. Lemma~3 gives both parts of Proposition~5.9 with the
source's condition $\tau\notin H^*$ retained for marked-IP transfer.
These are the source's explicitly borrowed results, not new results
claimed here. Their use in the small Table~2 attachment exclusions
still requires the exact local generator and support identifications.
Those classifications and the remaining full-source recovery are
outside this note. No software or numerical experiment is needed
for any assertion above.
\end{document}
